\section{Related Work And Positioning}

\subsection*{SAT-Based Computational Social Choice}

M4 uses a finite social-choice artifact in a tradition where SAT and automated
reasoning have been used to analyze impossibility and existence questions.
Computer-aided proofs of Arrow-style impossibility results, automated search
for impossibility theorems, and SAT-based synthesis or search for social
choice functions are well established \citep{tanglin2009computer,
geist2011automated,brandtgeist2016strategyproof}. Sen's liberal paradox is
the motivating social-choice setting for the present artifact
\citep{sen1970}. M4 is complementary to this line rather than a replacement
for it. It is not a new SAT search result, a speed result, or a new general
impossibility theorem. The finite Sen24 case is used here because it is small
enough to audit completely under the declared interface while still exposing
the gap between an impossibility report and a justified repair/reporting
claim.

This difference in emphasis matters. Much SAT-based social-choice work asks
whether a rule exists, whether an impossibility can be found, or whether a
small instance can illuminate a larger theorem. M4 assumes that finite
artifacts can already be useful and asks a later reporting question: once the
artifact reports inconsistency or repairability, which public claim is
licensed by the evidence actually recorded?

\subsection*{Certified SAT And Formal Verification}

The paper is also adjacent to certified SAT checking and proof-assistant work.
Lean~4 provides the proof-assistant setting for the semantic target used here
\citep{moura2021lean4}, while RAT/DRAT/LRAT-style checking belongs to the
broader ecosystem of proof-carrying SAT artifacts \citep{cruzfilipe2017lrat,
wetzler2014drat}. Mechanized social-choice formalization is an established
line of work in its own right \citep{nipkow2009social}. M4 draws on that
background by connecting its central right-condition target to a small
Lean-backed semantic definition. The claim remains limited: the Lean bridge
supports \code{RightAtomSemantics}, orientation invariance, and a sufficient
direction from two fixed right conditions to \code{MINLIB}. It does not prove
Python/CNF artifact correctness, semantic-to-CNF correctness, or full
selector-free equivalence to Lean \code{MINLIB}.

This places the paper between two stronger endpoints. It is more than an
uninterpreted finite table, because one central semantic target is tied to
Lean. It is less than an end-to-end certified artifact, because the generated
clauses and finite assignments are not proved equivalent to Lean social-choice
semantics. The point of the comparison is to make that middle position
explicit rather than to blur it.

\subsection*{Diagnosis, MUS/MCS, And Repair Explanation}

Repair and explanation questions connect M4 to diagnosis-oriented work such as
assumption-based truth maintenance and model-based diagnosis
\citep{dekleer1986atms,reiter1987diagnosis}. They also connect to modern
optimization and repair tooling, including MaxSAT-oriented views of minimal
relaxation \citep{bacchus2021maxsat}. M4 does not contribute a new diagnosis,
MUS/MCS, MaxSAT, or repair algorithm. Its concern is downstream of such
procedures: once an artifact has recorded conflicts, repairs, or grouped
reports, what exactly is the artifact entitled to say? The finite audit in
Section~\ref{sec:finite-audit} checks repair-empty SAT cells, UNSAT
orbit/fiber exactness, and support truncation for shape-blind reports under a
declared encoding. The methodological question is not how to find a repair
faster, but how to avoid over-reporting what the repair evidence licenses.

\subsection*{Abstraction, Reportability, And Claim Boundaries}

M4 is closest to work on explicit abstraction and reporting contracts.
Constraint hierarchies, for example, make tradeoffs among constraints explicit
rather than hiding them inside a single feasibility statement
\citep{borning1992constraint}. M4 uses a different but related discipline: it
separates finite audit evidence, declared semantic assumptions, a partial
Lean-backed semantic target, and future correctness obligations. Within the
Sen-Lean4 lineage, M1 made finite evidence diagnosable, M1.5 showed that raw
repair presentation need not be canonical, M2 established the semantic bridge
and general Sen theorem, and M3 supplied reportability conditions for grouped
repair reports. M4 does not re-prove those layers. It uses them to frame a
claim-boundary audit for the declared finite encoding.

The contribution is therefore methodological and scope-disciplined. M4 does
not solve AI governance or provide a general theory of automated decision
support. It gives a small, auditable example of how a formal artifact can say:
these finite claims are certified under the declared encoding; this semantic
target is Lean-backed; these stronger Python/CNF and semantic-to-CNF bridges
remain future correctness obligations.
That separation is the main novelty claimed here, not a new theorem of social
choice, a new SAT engine, or a new repair calculus.
